% SOURCE: Chinese PDF viewer pp.154--155, exactly repeated on pp.161--162.
\section{Particle-Number Representation}
The preceding method for treating identical-particle systems is first to assign a number to every particle, and then to consider all possible permutations of these numbers in order to satisfy exchange symmetry. This method has the following shortcomings: (1) since the particles are identical, they cannot in principle be distinguished by numbers; (2) when the number $N$ of particles is large, the symmetric or antisymmetric wavefunctions become very complicated; (3) the method cannot describe the conversion of particles into one another, and hence cannot be extended to a relativistic quantum many-particle system.

\subsection*{1) Fock space}
\textbf{Idea:} Use the eigenvalues and eigenstates of a single-particle observable $\op A$ to construct the eigenvalues and eigenstates of an identical-particle system, irrespective of whether the particles interact with one another.

The eigenvalues of $\op A$ are
\[a_1,a_2,\ldots;\]
the single-particle eigenstates are
\[\ket{a_1},\ket{a_2},\ldots;\]
and the particle numbers occupying these eigenstates are
\[n_1,n_2,\ldots.\]
According to the principle of indistinguishability, we do not know which particle occupies which particular single-particle eigenstate. We care only about the number of particles occupying each state. The state of the system is therefore described by
\[\ket{n_1,n_2,\ldots},\]
where $n_i$ is the number of particles occupying the state $\ket{a_i}$.

The multiparticle vector space formed by the states $\ket{n_1,n_2,\ldots}$ is called \emph{Fock space}. The state of an identical-particle system is described by the distribution of occupation numbers among the single-particle states. The vacuum and single-particle states can evidently be written
\[\ket 0=\ket{0,0,\ldots},\qquad
\ket{a_i}=\ket{0,0,\ldots,n_i=1,\ldots}.\]

\subsection*{2) Creation and annihilation operators}
Since a many-particle state is represented by its occupation-number distribution, the operators acting on this state change the occupation-number distribution. They are the particle creation and annihilation operators $\op a_i^\dagger$ and $\op a_i$:
\begin{align*}
\op a_i^\dagger\ket{n_1,\ldots,n_i,\ldots}&\sim\ket{n_1,\ldots,n_i+1,\ldots},\\
\op a_i\ket{n_1,\ldots,n_i,\ldots}&\sim\ket{n_1,\ldots,n_i-1,\ldots}.
\end{align*}
For the vacuum state,
\[\op a_i\ket0=0.\]
For a single-particle state,
\[\op a_i\ket{a_j}=\delta_{ij}\ket0,\qquad
\op a_i^\dagger\ket0=\ket{a_i}.\]

\subsection*{3) Transformation between two Fock spaces}
Consider two single-particle observables $\op A$ and $\op B$.
\[\begin{array}{c|cc}
&\text{$A$ representation}&\text{$B$ representation}\\\hline
\text{Eigenvalue}&a_i&b_i\\
\text{Eigenstate}&\ket{a_i}&\ket{b_i}\\
\text{Occupation number}&n_i&\widetilde n_i\\
\text{Creation operator}&\op a_i^\dagger&\op b_i^\dagger\\
\text{Annihilation operator}&\op a_i&\op b_i
\end{array}\]
The distribution of occupation numbers in the $\ket{a_i}$ states defines a Fock space $\ket{n_1,n_2,\ldots}$, while the distribution in the $\ket{b_i}$ states defines another Fock space $\ket{\widetilde n_1,\widetilde n_2,\ldots}$. The transformation between the two Fock spaces reduces to the transformation between $\op a_i,\op a_i^\dagger$ and $\op b_i,\op b_i^\dagger$.

The transformation of a single-particle eigenstate is
\[\ket{a_i}=\sum_j\ket{b_j}\braket{b_j}{a_i}
=\sum_jS_{ji}\ket{b_j},\qquad S_{ji}=\braket{b_j}{a_i}.\]
The representation-transformation matrix $S$ is unitary:
\[\sum_k S_{ki}^*S_{kj}=\delta_{ij},\qquad
\sum_k S_{ik}^*S_{jk}=\delta_{ij}.\]
Using the single-particle identity $\op a_i\ket{a_j}=\delta_{ij}\ket0$, the operator transformation gives
\begin{align*}
\op a_i\ket{a_j}
&=\delta_{ij}\ket0
=\sum_k S_{ki}^*S_{kj}\ket0
=\sum_{k,l}S_{ki}^*S_{lj}\delta_{kl}\ket0\\
&=\sum_{k,l}S_{ki}^*S_{lj}\op b_k\ket{b_l}
=\sum_kS_{ki}^*\op b_k\left(\sum_lS_{lj}\ket{b_l}\right)\\
&=\sum_kS_{ki}^*\op b_k\ket{a_j}.
\end{align*}
Thus,
\[\boxed{\op a_i=\sum_kS_{ki}^*\op b_k,\qquad
\op a_i^\dagger=\sum_kS_{ki}\op b_k^\dagger.}\]